\section{Conclusão}

Neste estudo foi realizada a comparação de diversos métodos de aprendizado de máquina e redes neurais artificiais para a detecção de patologias cardíacas em pacientes pediátricos. Os dados utilizados foram coletados pelo Real Hospital Português em colaboração com a Universidade do Porto. Os dados rotulados foram devidamente pré-processados e tratados antes de serem divididos em conjuntos de treino e teste. Para a otimização dos hiperparâmetros dos modelos, utilizou-se o método de busca em grade. A avaliação dos modelos foi realizada com base nas métricas de sensibilidade, especificidade e AUC.

Os resultados obtidos demonstram que a utilização de aprendizado de máquina e redes neurais artificias é efetiva para a detecção de patologias cardíacas em pacientes pediátricos. Além disso, a etapa de pré-processamento dos dados mostrou-se fundamental para a obtenção de bons resultados. A rede neural artificial \textit{Perceptron Multicamadas} obteve o melhor AUC (0,9490) entre os métodos avaliados, entretanto, o desempenho dos outros modelos foi muito próximo, com uma diferença máxima de 0,0191 em relação ao pior modelo. Isso evidencia que todos os modelos avaliados são eficazes para o problema, especialmente porque os dados foram adequadamente tratados e preparados para a modelagem.

Os resultados deste estudo são promissores e sugerem que o uso de aprendizado de máquina e redes neurais artificiais pode ser uma alternativa eficaz para a detecção de patologias cardíacas em pacientes pediátricos.

Pesquisas futuras podem explorar a combinação de múltiplas fontes de dados, o que pode aumentar a generalização dos modelos em diferentes hospitais e regiões. Além disso, é fundamental a validação dos modelos em um ambiente clínico real, para avaliar sua efetividade na prática.